\chapter*{Course Information}
\addcontentsline{toc}{chapter}{Course Information}
\section*{Course Objectives}
\begin{enumerate}[label=\arabic*.]
	\item To design transistor and operational-amplifier based electronic circuits.
	\item To simulate and implement the designed circuits in hardware and validate experimental results against theoretical/simulated predictions.
\end{enumerate}

\section*{Course Outcomes}
On successful completion of this course, the student will be able to:
\begin{table}[H]
	\centering
	\renewcommand{\arraystretch}{1.4}
	\begin{tabular}{cp{11cm}c}
		\toprule
		\textbf{CO} & \textbf{Course Outcome} & \textbf{Bloom's Level} \\
		\midrule
		CO1 & Use the various electronic instruments for conducting experiments. & K3 \\
		
		CO2 & Design and develop various electronic circuits using diodes and Zener diodes. & K3 \\
		
		CO3 & Design and implement amplifier and oscillator circuits using BJT and JFET. & K3 \\
		
		CO4 & Design and implement basic circuits using IC (OPAMP and 555 timer). & K3 \\
		
		CO5 & Simulate electronic circuits using circuit simulation software. & K3 \\
		
		CO6 & Use PCB layout software for circuit design. & K3 \\
		\bottomrule
	\end{tabular}
	\caption{Course Outcomes and Bloom's Knowledge Level}
\end{table}
\textit{K1: Remember, K2: Understand, K3: Apply, K4: Analyse, K5: Evaluate, K6: Create}

\section*{Scheme of Evaluation}

\subsection*{Continuous Internal Evaluation (CIE -- 50 Marks)}
\begin{table}[H]
	\centering
	\renewcommand{\arraystretch}{1.3}
	\begin{tabular}{p{8cm}c}
		\toprule
		\textbf{Component} & \textbf{Marks} \\
		\midrule
		Attendance & 5 \\
		
		Preparation / Pre-Lab Work, Viva, and Timely Completion of Lab Reports / Record (Continuous Assessment) & 25 \\
		
		Internal Examination & 20 \\
		
		\textbf{Total} & \textbf{50} \\
		\bottomrule
	\end{tabular}
	\caption{CIE Marks Distribution}
\end{table}

The continuous assessment component (25 marks) is further broken down as:
\begin{itemize}[label=\ding{226}]
	\item \textbf{Preparation and Pre-Lab Work (7 Marks):} Pre-lab assignments/quizzes and understanding of theoretical background.
	\item \textbf{Conduct of Experiments (7 Marks):} Adherence to procedure, accurate execution, safety, skill proficiency, and teamwork.
	\item \textbf{Lab Reports and Record Keeping (6 Marks):} Clarity, completeness, accuracy, and timely submission.
	\item \textbf{Viva Voce (5 Marks):} Ability to explain the experiment, results, and underlying principles.
\end{itemize}
The final marks for preparation, conduct, viva, and record are computed as the average across all experiments listed in the syllabus.

\subsection*{End Semester Examination (ESE -- 50 Marks)}
\begin{table}[H]
	\centering
	\renewcommand{\arraystretch}{1.3}
	\begin{tabular}{p{9.5cm}c}
		\toprule
		\textbf{Component} & \textbf{Marks} \\
		\midrule
		Procedure / Preliminary Work / Design / Algorithm & 10 \\
		
		Conduct of Experiment / Execution of Work / Troubleshooting / Programming & 15 \\
		
		Result with Valid Inference / Quality of Output & 10 \\
		
		Viva Voce & 10 \\
		
		Record & 5 \\
		\midrule
		\textbf{Total} & \textbf{50} \\
		\bottomrule
	\end{tabular}
	\caption{ESE Marks Distribution}
\end{table}

\begin{itemize}[label=\ding{226}]
	\item Students shall be permitted to appear for the end-semester examination only upon submission of a duly certified record.
	\item The external examiner shall endorse the record at the time of the examination.
\end{itemize}

\section*{CO--PO Mapping}
\begin{table}[H]
	\centering
	\renewcommand{\arraystretch}{1.3}
	\setlength{\tabcolsep}{4pt}
	\begin{tabular}{|c|c|c|c|c|c|c|c|c|c|c|c|c|}
		\hline
		& \textbf{PO1} & \textbf{PO2} & \textbf{PO3} & \textbf{PO4} & \textbf{PO5} & \textbf{PO6} & \textbf{PO7} & \textbf{PO8} & \textbf{PO9} & \textbf{PO10} & \textbf{PO11} & \textbf{PO12} \\
		\hline
		\textbf{CO1} & 3 & & & & & & & & & & & \\
		\hline
		\textbf{CO2} & 2 & 3 & 3 & 3 & 3 & & & & 3 & 3 & &  \\
		\hline
		\textbf{CO3} & 2 & 3 & 3 & 3 & 3 & & & & 3 & 3 & &  \\
		\hline
		\textbf{CO4} & 2 & 3 & 3 & 3 & 3 & & & & 3 & 3 & &  \\
		\hline
		\textbf{CO5} & 2 & 3 & 3 & 3 & 3 & & & & 3 & 3 & &  \\
		\hline
		\textbf{CO6} & 3 &   &   &   &   & & & & 3 & 3 & &  \\
		\hline
	\end{tabular}
	\caption{CO--PO Mapping (1: Slight, 2: Moderate, 3: Substantial, blank: No Correlation)}
\end{table}
%\textit{Note: PO columns reproduced from the syllabus mapping table as supplied; instructors should verify against the official curriculum document before final printing, since the source table's column alignment should be cross-checked.}

\section*{Text Books and References}
\begin{enumerate}[label=\arabic*.]
	\item Robert T.\ Paynter, \textit{Introductory Electronic Devices and Circuits}, Pearson Education.
	\item Boylestad R.\ L.\ and Nashelsky L., \textit{Electronic Devices and Circuit Theory}, Pearson Education.
	\item Donald A.\ Neamen, \textit{Electronic Circuits: Analysis and Design}, McGraw Hill.
\end{enumerate}
